\subsection{Blackwell 一维塌缩、冻结极端骨架与 InfoNCE 余维裂缝}\label{subsec:conclusion-blackwell-freezing-infonce-codimension-fracture}

沿用小节 \ref{subsec:conclusion-binfold-fixedn-chernoff-statdim}、小节 \ref{subsec:conclusion-freezing-ledger-blackwell-kolmogorov-likelihood-grid} 与小节 \ref{subsec:conclusion-capacity-majorization-infonce-complexity-separation} 的记号。前述结果已经分别给出：固定样本 escort 乘积试验在 Le Cam/Blackwell 意义下对一维二项实验的指数塌缩、冻结区中可逆恢复与 Bayes--InfoNCE 的双极账本分裂，以及完整 Bayes 最优 InfoNCE 数值塔对有序纤维谱的最小连续层析完备性。把这三条接口并读后，可再抽出一条更尖锐的维数分裂：固定样本的决策层在指数精度上只剩一维充分统计，而完整纤维热力学层析在 generic stratum 上仍需 \(|X_m|-1\) 个连续标量坐标。

\begin{theorem}[escort 决策层与热力层的余维裂缝]\label{thm:conclusion-escort-decision-thermal-codimension-gap}
\leanverified{paper\_conclusion\_escort\_decision\_thermal\_codimension\_gap}
固定 \(Q<\infty\) 与 \(n\in\NN\)。记
\[
\mathcal E_m^{(n,Q)}:=\{\pi_{m,q}^{\otimes n}:0\le q\le Q\},
\qquad
\mathcal G^{(n,Q)}:=\{\mathrm{Binomial}(n,\theta(q)):0\le q\le Q\},
\]
\[
\theta(q):=\frac{1}{1+\varphi^{q+1}},
\qquad
N_n(w^{(1)},\dots,w^{(n)}):=\sum_{j=1}^n w_m^{(j)}\in\{0,\dots,n\}.
\]
再记 \(N:=|X_m|\)，并令
\[
\Delta_N^\downarrow:=\bigl\{w\in[0,1]^N:\ w_1\ge\cdots\ge w_N,\ \sum_{i=1}^{N}w_i=1\bigr\},
\]
\[
U_N:=\bigl\{w\in\Delta_N^\downarrow:\ w_1>\cdots>w_N>0\bigr\}.
\]
则同时成立：
\begin{enumerate}
  \item fixed-sample escort 决策层在 Blackwell 意义下以指数速度塌缩到一维二项试验族 \(\mathcal G^{(n,Q)}\)，其前向充分统计量正是 \(N_n\)；
  \item 在 generic stratum \(U_N\) 上，任何由少于 \(N-1\) 个 Bayes 最优 InfoNCE 标量组成的连续读出，都不可能对完整有序纤维谱给出连续单射恢复，而完整族
  \[
  \{L_{K,m}^{\star}\}_{K=2}^{N}
  \]
  又足以唯一恢复之。
\end{enumerate}
因此，对固定样本数的 escort 判别，其有效统计维数为 \(1\)；而对完整纤维热力学层析，其最小连续恢复维数为 \(N-1\)。二者之间存在严格余维缺口
\[
(N-1)-1=N-2.
\]
\end{theorem}
\begin{proof}
第一条正是定理 \ref{thm:conclusion-binfold-escort-fixed-n-bounded-loss-decision-collapse} 的内容：任意 fixed-sample 有界损失决策规则都可在指数误差内替换为仅依赖 \(N_n\) 的规则，故 \(\mathcal E_m^{(n,Q)}\) 在 Le Cam/Blackwell 意义下塌缩到 \(\mathcal G^{(n,Q)}\)。

第二条由定理 \ref{thm:conclusion-capacity-ordered-spectrum-infonce-equivalence} 直接给出。该定理已说明：在 \(U_N\) 上，少于 \(N-1\) 个 Bayes 最优 InfoNCE 标量时不存在连续单射恢复，而完整数值塔 \(\{L_{K,m}^{\star}\}_{K=2}^{N}\) 与有序纤维谱严格等价。两条结论合并即得所述余维裂缝。证毕。
\end{proof}

\begin{theorem}[冻结区的极端骨架双坐标]\label{thm:conclusion-freezing-extremal-skeleton-two-coordinates}
\leanverified{paper\_conclusion\_freezing\_extremal\_skeleton\_two\_coordinates}
固定 \(q>0\) 并假设相应 escort 制度处于冻结区，记
\[
M_m:=\max_{x\in X_m} d_m(x),
\qquad
\nu_m:=\#\{x\in X_m:\ d_m(x)=M_m\}.
\]
若预算序列满足
\[
\frac{L_m}{M_m}\to \lambda\in[0,\infty),
\qquad
\frac{r_m-1}{\nu_m}\to \Lambda\in[0,\infty),
\]
则有
\[
\mathrm{Succ}^{\mathrm{rev}}_{m,q}(L_m)\longrightarrow \min(1,\lambda),
\qquad
\mathcal L^\star_{m,q}(r_m)\longrightarrow
\EE\bigl[\log(1+\mathrm{Poisson}(\Lambda))\bigr].
\]
因此，冻结区主导可操作信息的 leading skeleton 严格塌缩为
\[
\mathfrak E_m:=(M_m,\nu_m).
\]
换言之，可逆恢复极限只记住最大纤维尺度 \(M_m\)，而 Bayes--InfoNCE 极限只记住最大层简并度 \(\nu_m\)。
\end{theorem}
\begin{proof}
这正是定理 \ref{thm:conclusion-freezing-dual-ledger-completeness} 在单一系统上的直接特例。前者由 \(L_m/M_m\) 的极限控制，后者由 \((r_m-1)/\nu_m\) 的极限控制，且二者分别收敛到题设中的闭式极限。故冻结区的 leading operational data 恰由 \((M_m,\nu_m)\) 封闭。证毕。
\end{proof}

\begin{corollary}[全谱完备性与极端塌缩的并存]\label{cor:conclusion-full-spectrum-completeness-coexists-with-extremal-collapse}
\leanverified{paper\_conclusion\_full\_spectrum\_completeness\_coexists\_with\_extremal\_collapse}
固定分辨率 \(m\) 时，完整 Bayes 最优 InfoNCE 数值塔
\[
\{L_{K,m}^{\star}\}_{K=2}^{|X_m|}
\]
唯一决定完整有序纤维谱；但在冻结缩放极限中，同一对比学习体系的主导极限只记住 \(\nu_m\)，而最优可逆恢复门槛只记住 \(M_m\)。因此，有限层析意义下的全谱完备性与极限操作意义下的极端骨架塌缩并不矛盾，而是同一可观测体系在两个层次上的严格分裂。
\end{corollary}
\begin{proof}
由定理 \ref{thm:conclusion-capacity-ordered-spectrum-infonce-equivalence}，完整数值塔 \(\{L_{K,m}^{\star}\}_{K=2}^{|X_m|}\) 与有序纤维谱严格等价，故在固定 \(m\) 上给出全谱完备层析。另一方面，定理 \ref{thm:conclusion-freezing-extremal-skeleton-two-coordinates} 已说明冻结区中的可逆恢复与 Bayes--InfoNCE leading limit 分别只由 \(M_m\) 与 \(\nu_m\) 控制。两者对应的正是 finite tomography 与 asymptotic operation 两个不同层次，故可以严格并存。证毕。
\end{proof}

\paragraph{统一读法}
这条链把 escort 统计、冻结操作与 Bayes 对比层析压成一个清晰的双层结构：固定样本数下的全部有界损失判别在指数精度内退化为末位计数 \(N_n\) 上的一维二项实验；而一旦要求连续恢复完整有序纤维谱，generic stratum 上又必须支付 \(|X_m|-1\) 个标量坐标的最小维数。与此同时，冻结极限再把操作层压回两枚极端参数 \(M_m\) 与 \(\nu_m\)。因此，本文中的统计实验、热力层析与极限操作并不共享同一“有效维数”，而是严格分裂为一维决策骨架、二维极端骨架与 \((|X_m|-1)\)-维全谱坐标三种互不塌缩的资源层。

\endinput
